\documentclass[10pt,twocolumn]{article}

\usepackage[margin=0.75in]{geometry}
\usepackage{graphicx}
\usepackage{amsmath, amssymb}
\usepackage{algorithm}
\usepackage{algpseudocode}
\usepackage{booktabs}
\usepackage{url}
\usepackage{hyperref}
\usepackage{authblk}
\usepackage{caption}

\setlength{\columnsep}{0.25in}
\setcounter{secnumdepth}{3}

\title{\vspace{-2em}Mitigating Bias in Facial Recognition Systems:\\
A Corrected Replication of DB-VAE on UTKFace and Tiny ImageNet}

\author{Kwadwo Twumasi \\ CUNY Lehman College \\ Mentor: Professor Liang Zhao}
\date{}

\begin{document}
\maketitle
\begin{abstract}
{\normalsize
\noindent
As machine learning-based systems become more prevalent, it is important to examine methods to mitigate bias within these systems, especially bias against underrepresented segments of society in training data.
In this research project, our goal is to evaluate the generalization and applicability of the algorithmic bias mitigation approach designed by MIT researchers (Amini et al., 2019), which combines a variational autoencoder (VAE) with a binary classification task to learn latent features and adaptively re-weight training examples by their rarity in the learned latent space, controlled by a debiasing hyperparameter $\alpha$.
To examine how well the Debiasing Variational Autoencoder (DB-VAE) algorithm generalizes across dataset compositions, we replicate and extend this method on a different dataset combination: UTKFace~\cite{zhang2017utkface} (positives) and Tiny ImageNet~\cite{le2015tinyimagenet} (negatives), evaluated across a 10-group race~$\times$~gender held-out split.
In doing so, we identify and correct a loss-function bug in the widely used MIT 6.S191 Lab~2 template that swaps the face-indicator factor between the classification and VAE terms, zeroing out the classification gradient on non-face training samples.
We then conduct a comprehensive $\alpha$-sweep over $\{0.001, 0.01, 0.05, 0.1, \infty\}$ across three random seeds.
Despite the corrected formulation, the baseline classifier achieves $\mathbb{E}[\mathcal{A}] = 0.9982$ across all ten groups, and the $\alpha$-sweep produces no measurable change in either mean accuracy or cross-group variance.
We interpret this finding as a characterization of the regime in which latent-density resampling has headroom to act and in which it does not.
This suggests that a class-balanced training pool of cleanly-cropped faces leaves minimal opportunity for the debiasing intervention to improve upon the baseline.
}
    \end{abstract}

% =====================================================================
\section{Introduction}
% =====================================================================

Machine learning systems are increasingly used in consequential
decisions~\cite{klare2012face}, and a recurring failure mode is that they
perform worse on demographic groups underrepresented in their training
data. Facial detection is one of the most studied instances of this
failure: commercial systems have been shown to misclassify darker-skinned
women at substantially higher rates than lighter-skinned
men~\cite{buolamwini2018gendershades}, and academic benchmarks now exist
specifically to surface such disparities.

A natural mitigation is to up-weight training examples whose features are
rare in the training set, so that the gradient signal during training is
spread more evenly across the latent space of the positive class.
The Debiasing Variational Autoencoder of Amini et
al.~\cite{amini2019dbvae} (DB-VAE) implements this idea end-to-end: a
single network simultaneously learns the binary classification task and an
unsupervised variational latent representation of the positive class, and
the inverse density of each training example in that latent space is used
as its sample weight at the next epoch.
The strength of the intervention is
controlled by a single smoothing hyperparameter~$\alpha$: small~$\alpha$
pushes the sampled distribution toward uniformity over the latent space,
while~$\alpha \to \infty$ reverts to uniform sampling over the original
training set.

The original DB-VAE paper demonstrates the method on CelebA~\cite{liu2015celeba} and ImageNet~\cite{deng2009imagenet},
evaluated on the Pilot Parliaments Benchmark (PPB)~\cite{buolamwini2018gendershades}
across four demographic bins (lighter/darker skin, male/female), and
reports that decreasing~$\alpha$ both increases the mean accuracy and decreases
its cross-group variance. This work asks a complementary question: does the
DB-VAE mechanism generalize to a different dataset combination, and how
does it behave on a finer-grained demographic split?

We make the following contributions:
\begin{enumerate}
    \item \textbf{Methodological.} We identify a loss-function bug in the
    MIT 6.S191 Lab~2 PyTorch template, widely used as the starting point
    for DB-VAE re-implementations. The bug swaps a face-indicator factor
    between the classification and VAE terms, zeroing out classification
    loss on non-face samples and degenerating the trained model into a
    yes-face oracle. We give a one-line correction and a diagnostic
    procedure that detects failure in a single forward pass.
    \item \textbf{Empirical.} We replicate the DB-VAE architecture and
    training procedure on UTKFace~\cite{zhang2017utkface} as the positive
    class and Tiny ImageNet~\cite{le2015tinyimagenet} as the negative
    class, with a 10-group race~$\times$~gender evaluation split
    (5~UTKFace race categories~$\times$~2 genders, 200 held-out faces per
    group) that subsumes the standard 4-group skin-tone~$\times$~gender
    scheme.
    \item \textbf{Negative result.} On this dataset combination, the
    corrected baseline classifier achieves $\mathbb{E}[\mathcal{A}] =
    0.9982$ across all ten groups, and an $\alpha$-sweep matching the
    grid used by Amini et al.\ produces no measurable improvement over
    the no-debiasing baseline. We argue this outcome is itself diagnostic:
    the DB-VAE intervention can only act on headroom that exists in the
    baseline, and a class-balanced training pool of cleanly-cropped faces
    leaves essentially none.
\end{enumerate}

% =====================================================================
\section{Related Work}
% =====================================================================

DB-VAE belongs to the in-processing group of fairness interventions, which operate during model training rather than modifying the data beforehand or post-processing model outputs.

\paragraph{Resampling for class imbalance.}
Naive oversampling of a minority class~\cite{more2016survey} addresses imbalance \emph{across} classes but is silent on imbalance \emph{within} a class -- the under-representation of, say, darker-skinned faces inside a broader ``face'' category. DB-VAE addresses within-class imbalance by learning the latent structure of the class and inverting that distribution at sample time.

\paragraph{Variational autoencoders.}
The encoder--decoder structure DB-VAE inherits is the standard variational autoencoder of Kingma and Welling~\cite{kingma2013vae} and Rezende et al.~\cite{rezende2014stochastic}. The novelty in DB-VAE is in how the learned latent variables are used at training time, not in the architecture itself.

\paragraph{Fairness benchmarks for face systems.}
PPB~\cite{buolamwini2018gendershades} is the canonical benchmark for intersectional facial-analysis bias and uses Fitzpatrick skin type as a proxy for race.
UTKFace~\cite{zhang2017utkface} provides explicit race and gender labels across five race categories and binary gender, making it particularly well-suited for examining the model's race and gender debiasing capabilities.
We selected UTKFace specifically because its structured annotations allow us to evaluate performance across 10 intersectional race~$\times$~gender groups, providing finer-grained demographic analysis than the four-bin scheme in the original DB-VAE paper.
The choice of benchmark substantially affects the difficulty of the underlying task, as we will return to in Section 5.

% =====================================================================
\section{Methodology}
% =====================================================================

\subsection{DB-VAE architecture and loss}
\label{sec:arch}

The DB-VAE shares an encoder with the baseline classifier and adds a
decoder that mirrors it. For binary classification with~$k$ latent
variables, the encoder produces~$2k+1$ activations: one classification
logit~$\hat{y}$, $k$ latent means~$\boldsymbol{\mu}$, and~$k$ latent
log-standard-deviations~$\log\boldsymbol{\sigma}$. A latent
sample~$\mathbf{z}$ is drawn via the reparameterization trick
$\mathbf{z} = \boldsymbol{\mu} + \exp(\log\boldsymbol{\sigma}) \odot
\boldsymbol{\epsilon}$ with~$\boldsymbol{\epsilon} \sim \mathcal{N}(0, I)$,
and a decoder reconstructs~$\hat{\mathbf{x}}$ from~$\mathbf{z}$.

The training loss combines classification, reconstruction, and a KL
regularizer:
\begin{equation}
\begin{aligned}
    \mathcal{L}_{\text{TOTAL}} = {} & c_1 \underbrace{\left[ \sum_{i \in \{0, 1\}} y_i \log \left( \frac{1}{\hat{y}_i} \right) \right]}_{\mathcal{L}_y(y, \hat{y})} + c_2 \underbrace{\left[ \|\mathbf{x} - \hat{\mathbf{x}}\|_p \right]}_{\mathcal{L}_x(\mathbf{x}, \hat{\mathbf{x}})} \\
                                     & + c_3 \underbrace{\left[ \frac{1}{2} \sum_{j=0}^{k-1} \left( \sigma_j + \mu_j^2 - 1 - \log(\sigma_j) \right) \right]}_{\mathcal{L}_{KL}(\boldsymbol{\mu}, \boldsymbol{\sigma})},
    \label{eq:loss}
\end{aligned}
\end{equation}
where $c_1$, $c_2$, and $c_3$ weight the classification, reconstruction, and KL terms, respectively; in our experiments we set~$c_1 = 1$ and~$c_2 = c_3 = 5 \times 10^{-4}$. The encoder backbone matches the
baseline CNN (four convolutional blocks of 12, 24, 48, 72 filters with
batch normalization~\cite{ioffe2015batchnorm}, then a 512-unit dense layer);
the decoder is its mirror image with transposed convolutions. We
use~$k = 100$ latent variables, matching Amini et al.~\cite{amini2019dbvae}.

\subsection{A loss-function bug in the MIT lab template}
\label{sec:bug}

The PyTorch implementation of Equation (1) distributed in the MIT
6.S191 Lab~2 template -- the most accessible reference implementation of
DB-VAE -- combines its components as
\begin{equation}
    \mathcal{L}_{\text{MIT}} = \mathrm{mean}\bigl(\mathcal{L}_y \cdot \mathbb{I}_f + \mathcal{L}_{\text{VAE}}\bigr),
    \label{eq:buggy}
\end{equation}
swapping the face-indicator factor between the two terms relative to
Equation (1). Two consequences follow. First, the classification loss
on non-face examples is multiplied by zero, so the gradient with respect to
$\hat{y}$ provides no signal to push down probabilities on non-faces.
Second, the VAE reconstruction and KL losses are evaluated on non-face
examples (a Tiny ImageNet animal, a vehicle, a piece of furniture), which
the original method explicitly prohibits.

The combination is sufficient to collapse the trained model into a constant
predictor that outputs $\sigma(\hat{y}) \approx 1$ on every input. The
collapse is invisible to any evaluation that uses only face inputs (which
all standard DB-VAE evaluations do, including the original paper's PPB
protocol), so the failure can persist undetected across many runs.

We detect the bug by evaluating the trained model on a balanced batch of
500 faces and 500 non-faces drawn from the same training loader. Before
the fix, P(face) on non-faces was $1.0000 \pm 0.0000$ (mean over the
non-face half of the batch) and non-face accuracy was $0.0000$. After
applying the correction of Equation (1), P(face) on non-faces drops
to $0.0002$ (max $0.0332$) and non-face accuracy rises to $1.0000$ -- a
five-thousand-fold change in the diagnostic statistic. All experimental
results in Section 4.2 use the corrected formulation.

\subsection{Adaptive resampling}
\label{sec:resample}

At the start of each epoch, the encoder is run over all training faces, and
the marginal distribution~$\hat{Q}_i(z_i | X)$ of each latent dimension is
approximated by a 10-bin histogram. The training sampling weight for
face~$\mathbf{x}$ is
\begin{equation}
    \mathcal{W}(\mathbf{z}(\mathbf{x}) | X) \;\propto\;
    \prod_{i}\;\frac{1}{\hat{Q}_i(z_i(\mathbf{x}) | X) + \alpha},
    \label{eq:weights}
\end{equation}
where the smoothing constant~$\alpha$ controls how aggressively the
resampling pushes toward uniformity over the latent space.
As~$\alpha \to 0$ the sampled distribution approaches uniform over the
latent space; as~$\alpha \to \infty$ it approaches uniform sampling over
the original training set. We follow Amini et al. in taking the
element-wise maximum across latent dimensions, which avoids vanishing
weights when any single dimension exhibits a sharp density spike.
Algorithm 1 summarizes the procedure.

\begin{algorithm}[H]
    \caption{DB-VAE training with adaptive resampling}
    \label{alg:dbvae}
    \begin{algorithmic}[1]
        \Require Training set $\{X, Y\}$, batch size $b$, debiasing strength $\alpha$
        \State Initialize encoder, decoder weights
        \For{each epoch}
            \State Compute $\boldsymbol{\mu}_i$ for every face $\mathbf{x}_i \in X$
            \State Update histograms $\hat{Q}_i(z_i | X)$
            \State $\mathcal{W}(\mathbf{z}(\mathbf{x}) | X) \gets \prod_{i} \dfrac{1}{\hat{Q}_i(z_i(\mathbf{x}) | X) + \alpha}$
            \For{each batch}
                \State Sample faces $\sim \mathcal{W}$, non-faces uniformly
                \State Forward, compute $\mathcal{L}$ via Equation (1), backpropagate
                \State Update encoder and decoder parameters
            \EndFor
        \EndFor
    \end{algorithmic}
\end{algorithm}

\subsection{Ten-group evaluation protocol}
\label{sec:eval}

Amini et al.~\cite{amini2019dbvae} report performance over four
intersectional groups using Fitzpatrick skin type as a race proxy. UTKFace
ships with five explicit race labels (White, Black, Asian, Indian, Other)
and a binary gender label, so we report ten intersectional groups:
\textit{\{race\}~$\times$~\{Female, Male\}}. Two hundred faces per group
are held out before training and never seen during the latent-distribution
or weight-update steps. For each group, we report the accuracy at threshold
0.5 (the fraction of test faces for which $\sigma(\hat{y}) > 0.5$). The
aggregate metrics follow Amini et al.'s Table~1: $\mathbb{E}[\mathcal{A}]$
is the mean of the per-group accuracies and $\mathrm{Var}[\mathcal{A}]$
their population variance.

This protocol differs from the original DB-VAE evaluation in two
methodologically important ways. First, our test images are pre-cropped
faces rather than wide-angle parliamentarian photos, so the task is one of
classification rather than sliding-window detection. Second, our ten-group
split preserves race information that the four-group skin-tone scheme
collapses, allowing us to ask whether DB-VAE's effect is uniform across
race categories within a single skin-tone bin or whether it varies. The
practical consequence of the first difference -- a substantially easier
underlying task -- is the central observation of Section 4.2.

% =====================================================================
\section{Experiments}
% =====================================================================

\subsection{Datasets and training setup}

We use UTKFace~\cite{zhang2017utkface} as the positive class: 21{,}705 aligned face crops after
removing the held-out evaluation set, resized to $64 \times 64$ pixels.
The negative class is a 25{,}000-image random subsample of Tiny ImageNet~\cite{le2015tinyimagenet}
(125 images per class across 200 generic object classes), also at
$64 \times 64$. The aggregate training pool is therefore 46{,}705 images,
roughly balanced between faces and non-faces, contrasting with the
    $400{,}000$-image, CelebA~\cite{liu2015celeba}-derived training set used by Amini et al.

All models train with Adam (lr~$=5 \times 10^{-4}$), batch size~32, for ten
epochs. The DB-VAE uses 100 latent dimensions and a KL weight
of~$5 \times 10^{-4}$. We sweep~$\alpha \in \{\infty, 0.1, 0.05, 0.01,
0.001\}$, where~$\alpha = \infty$ denotes the no-debiasing baseline in
which the resampling step is skipped, and the DB-VAE is trained with uniform
face sampling. Each~$\alpha$ is run with three random seeds (initial
weights, batch order, reparameterization noise), and we report the
mean~$\pm$~SEM (standard error of the mean) across seeds. The full sweep of 15 runs takes approximately
57~minutes on a single RTX 3060 Laptop GPU.

\subsection{Per-group accuracy}
\label{sec:results}

\begin{figure*}[t]
    \centering
    \includegraphics[width=0.92\textwidth]{c:/Users/Savag/OneDrive/Desktop/capstone_project/report/figures/per_group_accuracy.png}
    \caption{Per-group accuracy at threshold 0.5 across the ten
    race~$\times$~gender groups, for the no-debiasing baseline and four
    DB-VAE conditions. Bars show the mean across three seeds; error bars are
    one SEM. The y-axis is zoomed to the relevant range to show differences across conditions.}
    \label{fig:per-group}
\end{figure*}

\begin{figure*}[t]
    \centering
    \includegraphics[width=0.92\textwidth]{c:/Users/Savag/OneDrive/Desktop/capstone_project/report/figures/aggregate_metrics.png}
    \caption{Aggregate metrics across the $\alpha$-sweep. Mean accuracy $\mathbb{E}[\mathcal{A}]$ (left) and cross-group variance $\mathrm{Var}[\mathcal{A}]$ (right) are plotted as a function of debiasing strength $\alpha$. Error bars represent one SEM across three random seeds. Despite the corrected loss formulation, the $\alpha$-sweep produces no measurable improvement in either metric over the no-debiasing baseline.}
    \label{fig:aggregate}
\end{figure*}

\begin{table}[H]
    \centering
    \small
    \caption{Aggregate accuracy and bias on the ten-group held-out test
    set after applying the loss-function correction of
    Section 3.2. $\mathbb{E}[\mathcal{A}]$ is the mean and
    $\mathrm{Var}[\mathcal{A}]$ the population variance of the per-group
    accuracies. Values are the mean~$\pm$~SEM over three seeds.}
    \label{tab:agg}
    \begin{tabular}{lcc}
        \toprule
        Condition & $\mathbb{E}[\mathcal{A}]$ & $\mathrm{Var}[\mathcal{A}]\ (\times 10^{-5})$ \\
        \midrule
        No debias       & $0.9982 \pm 0.0004$ & $0.875 \pm 0.238$ \\
        $\alpha = 0.1$  & $0.9980 \pm 0.0003$ & $1.083 \pm 0.030$ \\
        $\alpha = 0.05$ & $0.9970 \pm 0.0005$ & $1.217 \pm 0.156$ \\
        $\alpha = 0.01$ & $0.9977 \pm 0.0002$ & $0.950 \pm 0.163$ \\
        $\alpha = 0.001$& $0.9975 \pm 0.0003$ & $1.275 \pm 0.163$ \\
        \bottomrule
    \end{tabular}
\end{table}

Figure 1 shows per-group accuracies for
each~$\alpha$ condition, with the y-axis zoomed to the actual range of
the data. Table 1 summarizes the aggregate
$\mathbb{E}[\mathcal{A}]$ and $\mathrm{Var}[\mathcal{A}]$. Three
observations stand out.

First, the corrected baseline classifier achieves a mean accuracy
$\mathbb{E}[\mathcal{A}] = 0.9982$ across all ten demographic groups.
Eight of the ten groups sit at $\geq 99.5\%$ in the no-debias condition,
the two consistent under-performers are White Male (99.2\%) and Asian Male
(99.5\%). There is no demographic group on which the baseline exhibits the
20--30 percentage-point gap that motivated the DB-VAE method on PPB.

Second, the $\alpha$-sweep produces no monotonic trend in either
$\mathbb{E}[\mathcal{A}]$ or $\mathrm{Var}[\mathcal{A}]$. The differences
across $\alpha$ values lie within or close to one SEM of the no-debiasing
baseline. The no-debiasing condition has the highest
$\mathbb{E}[\mathcal{A}]$ (0.9982) and the lowest
$\mathrm{Var}[\mathcal{A}]$ ($8.75 \times 10^{-6}$) of any condition,
though not by a margin that survives the seed-to-seed noise floor. For
comparison, the Amini et al.\ Table~1 shows $\mathbb{E}[\mathcal{A}]$
moving from 95.13\% to 97.36\% and $\mathrm{Var}[\mathcal{A}]$ from 28.84
to 9.43 (in percentage-point units squared) over the same $\alpha$ grid --
changes orders of magnitude larger than what we observe.

Third, the per-group SEMs are themselves on the order of
$1.7 \times 10^{-3}$. The within-condition group-to-group differences --
say, the 99.2\% vs.~99.5\% gap between White Male and Asian Male -- are at
or below this noise floor. Whatever bias the corrected DB-VAE could, in
principle, redress on this benchmark, the metric does not give it room to
move.

\subsection{Why the bug was invisible to standard evaluation}

It is worth noting that the loss-function bug of Section 3.2
would have been undetectable using only the DB-VAE-vs-Standard-CNN
comparison plot included in the MIT lab template. That comparison evaluates
each model on a face-only test set (the 10 race~$\times$~gender groups), and
a constant ``yes-face'' predictor scores 100\% on a face-only test by
construction. The diagnostic we propose (Section 3.2) instead
asks the trained model what probability it assigns to mixed face and
non-face inputs; this single forward pass reveals the collapse in one
number. We recommend that any DB-VAE replication include this check.

% =====================================================================
\section{Discussion}
\label{sec:discussion}
% =====================================================================

\subsection{Two interpretations of the saturated null result}

The corrected DB-VAE classifier produces a baseline that is already at
$\geq 99.2\%$ on every demographic group, and the resampling
intervention (Equation (2)) does not change this. There are two
non-exclusive interpretations.

\emph{(1) The task is too easy.} The positive class consists of aligned
UTKFace crops -- centered, frontal, well-lit -- whereas the original
DB-VAE experiments evaluated on the Pilot Parliaments Benchmark using a
sliding-window detection protocol, which introduces variability in pose,
illumination, and crop framing that the cropped-classification task
removes. The negative class, Tiny ImageNet, is dissimilar enough from
faces that the binary boundary is easy to learn. Either change
individually narrows the gap between baseline accuracy and the ceiling;
together, they essentially eliminate it.

\emph{(2) The training set is already balanced.} UTKFace's five race
labels and two genders give a more uniform demographic distribution than
    the heavily light-skinned CelebA~\cite{liu2015celeba} used by Amini et al.\ -- with the result
that the latent-density variation that DB-VAE re-weights against is
itself smaller. The largest under-representation in our training set is
about $2 \times$ between the most and least populous race category; the
analogous CelebA~\cite{liu2015celeba} imbalance for dark-skinned subjects is closer to
$10 \times$.

Either way, Equation (2) re-weights training examples toward
those whose latent codes are rare. If the classifier already generalizes
correctly to every demographic group, increasing the representation of any
subgroup leaves the gradient signal unchanged at convergence: there is no
group whose loss the resampling could reduce.

\subsection{Implications for practice}

We see the result as a positive contribution to the conditions under
which DB-VAE is useful, not a refutation of the method. DB-VAE was
designed for and validated on a setting where the baseline exhibits
substantial, demographically structured failures (CelebA~\cite{liu2015celeba} evaluated on
PPB). When practitioners adopt the technique, they should first verify
that such failures exist on their benchmark; otherwise, the resampling
adds training cost without delivering the benefit it was designed to
provide.

The bug we describe in Section 3.2 is more directly actionable.
The MIT 6.S191 Lab~2 template is widely used as an entry point for
researchers learning DB-VAE, and the failure mode it produces (a constant
yes-face predictor) is invisible to standard face-only evaluation. Any
replication of DB-VAE that begins from this template should apply the
one-line correction and run the mixed-input diagnostic before drawing
conclusions about $\alpha$.

\subsection{Limitations}

\textit{(i)}~Our training pool is roughly an order of magnitude smaller
than the $400{,}000$-image set used by Amini et al., which itself bounds
how extreme the latent-density variation across groups can be.
\textit{(ii)}~We did not evaluate under test-time perturbation (noise,
blur, partial occlusion) that would introduce headroom artificially; this
is a natural follow-up experiment.
\textit{(iii)}~We did not deliberately bias the training subsample
(e.g., by subsampling UTKFace down to a 70\%-White-Male composition),
which would be the cleanest way to recover the regime DB-VAE was designed
for.
\textit{(iv)}~Three seeds is fewer than the five used in the original
paper; we verified post-hoc that the seed-to-seed SEM is small enough
that this choice does not change the qualitative finding.

% =====================================================================
\section{Conclusion}
% =====================================================================

We replicated the DB-VAE adaptive-resampling debiasing method on a new
dataset combination (UTKFace and Tiny ImageNet) and a finer-grained
ten-group race~$\times$~gender evaluation. Two findings stand out.

First, the MIT 6.S191 Lab~2 PyTorch template contains a loss-function bug
that collapses the trained DB-VAE into a constant yes-face predictor; the
failure is invisible to the standard face-only evaluation and the
template's own visualization cells. We give a one-line correction matching
Amini et al.\ Sec.~3.2 and a single-pass diagnostic for detecting the
collapse.

Second, with the corrected loss, the baseline classifier on UTKFace and
Tiny ImageNet already achieves 99.8\% mean accuracy across all ten
demographic groups, and the $\alpha$-sweep that was decisive in the
original paper produces no measurable effect. We interpret this as a
characterization of DB-VAE's regime of applicability: the method requires
latent-density variation across groups that the baseline cannot already
absorb. Future work should re-evaluate on a deliberately biased training
subsample of UTKFace or a perturbed test-time protocol -- conditions under which
the saturated-baseline regime no longer holds. The full pipeline,
including the local Kaggle-driven data loaders, the bug diagnostic, and
the corrected sweep harness, is implemented in
PyTorch~\cite{paszke2019pytorch} and runs end-to-end on a single
consumer-grade GPU.

% =====================================================================
\section*{Acknowledgments}
% =====================================================================

This report was drafted with assistance from AI tools (Grammerly and GitHub Copilot CLI) for writing refinement and formulating the initial LaTeX document code. We gratefully acknowledge Professor Liang Zhao for guiding dataset selection, supervising the experimental design, and providing valuable feedback on the analysis of results. All experimental work, analysis, and research contributions are the author's own.

{\small
\bibliographystyle{plain}
\bibliography{references}
}

\end{document}